\documentclass{article}
\usepackage{fontspec}
\usepackage{tikz}
\usepackage{tikz-dependency}
\usepackage{adjustbox}


% === Your polytonic Greek setup (customize with your purchased fonts) ===
\defaultfontfeatures{Ligatures=TeX}
\setmainfont{ArnoPro-Regular}          % e.g. Brill, IFAO, or your professional font
% \newfontfamily\greekfont[Script=Greek]{YourGreekFontHere}
\newfontfamily\greekfont[Script=Greek]{Arno Pro}

% \usepackage{polyglossia}
% \setotherlanguage[variant=ancient]{greek}


\begin{document}
\begin{figure}[ht]
\centering
\begin{adjustbox}{max width=\textwidth}
\begin{tikzpicture}
  \begin{deptext}[column sep=0.9em]
    τῇ \& βασιλέως \& θυγατρὶ \& ἐδήλου \& ὁ \& ποιητὴς \& τὰ \& ποιήματα \& τὰ \& περὶ \& τῆς \& φύσεως \& γεγραμμένα \\
  \end{deptext}
  \draw[blue!70!25, fill=blue!70!12, rounded corners=3pt, line width=0.8pt] ([xshift=-0.3em]deptext-1-1.north west) rectangle ([xshift=0.3em]deptext-1-13.south east);
  \draw[red!70!25, fill=red!70!12, rounded corners=3pt, line width=0.8pt] ([xshift=-0.3em]deptext-1-7.north west) rectangle ([xshift=0.3em]deptext-1-13.south east);
  \node[draw, fill=blue!70!20, rounded corners=2pt, inner sep=2pt, font=\small] at (deptext-1-1) {τῇ};
  \node[draw, fill=blue!70!20, rounded corners=2pt, inner sep=2pt, font=\small] at (deptext-1-2) {βασιλέως};
  \node[draw, fill=blue!70!20, rounded corners=2pt, inner sep=2pt, font=\small] at (deptext-1-3) {θυγατρὶ};
  \node[draw, fill=blue!70!20, rounded corners=2pt, inner sep=2pt, font=\small] at (deptext-1-4) {ἐδήλου};
  \node[draw, fill=blue!70!20, rounded corners=2pt, inner sep=2pt, font=\small] at (deptext-1-5) {ὁ};
  \node[draw, fill=blue!70!20, rounded corners=2pt, inner sep=2pt, font=\small] at (deptext-1-6) {ποιητὴς};
  \node[draw, fill=blue!70!20, rounded corners=2pt, inner sep=2pt, font=\small, draw=red!70, line width=1.2pt] at (deptext-1-7) {τὰ};
  \node[draw, fill=blue!70!20, rounded corners=2pt, inner sep=2pt, font=\small, draw=red!70, line width=1.2pt] at (deptext-1-8) {ποιήματα};
  \node[draw, fill=blue!70!20, rounded corners=2pt, inner sep=2pt, font=\small, draw=red!70, line width=1.2pt] at (deptext-1-9) {τὰ};
  \node[draw, fill=blue!70!20, rounded corners=2pt, inner sep=2pt, font=\small, draw=red!70, line width=1.2pt] at (deptext-1-10) {περὶ};
  \node[draw, fill=blue!70!20, rounded corners=2pt, inner sep=2pt, font=\small, draw=red!70, line width=1.2pt] at (deptext-1-11) {τῆς};
  \node[draw, fill=blue!70!20, rounded corners=2pt, inner sep=2pt, font=\small, draw=red!70, line width=1.2pt] at (deptext-1-12) {φύσεως};
  \node[draw, fill=blue!70!20, rounded corners=2pt, inner sep=2pt, font=\small, draw=red!70, line width=1.2pt] at (deptext-1-13) {γεγραμμένα};
\end{tikzpicture}

\vspace{0.4em}
\begin{tabular}{@{}ll@{}}
  \colorbox{blue!70!25}{\rule{0.9cm}{0.45cm}} & \textbf{VU1} (Level 1): Independent Clause — Transitive \\
  \colorbox{red!70!25}{\rule{0.9cm}{0.45cm}} & \textbf{VU2} (Level 2): Attributive Participle — Intransitive \\
\end{tabular}
\end{adjustbox}

\caption{Verbal Units — Christopher Blackwell — τῇ βασιλέως θυγατρὶ ἐδήλου ὁ ποιητὴς τὰ ποιήματα τὰ περὶ τῆς φύσεως γεγραμμένα .}
\end{figure}
\end{document}
